\section{The strongest objections}\label{sec:objections}
\begin{enumerate}
\item \textbf{CBT has a large evidence base.} Yes; Section~\ref{sec:evidence}. The argument is about mechanism and moderators, not average efficacy.
\item \textbf{CBT for adult ADHD works.} A trial shows benefit over a matched control \citep{safren2010cbt}, but in medicated adults, with a skills-heavy protocol, and the Cochrane review grades the evidence as low quality \citep{lopez2018cochrane}.
\item \textbf{Patients practice when calm.} Accepted in Section~\ref{sec:ingredients}; it moves the active ingredient to practice.
\item \textbf{Exposure is behavioural.} Likewise; it strengthens the argument against restructuring specifically.
\item \textbf{Mindfulness is not better.} True (Section~\ref{sec:alternatives}); the argument does not require superiority, only that restructuring is not uniquely necessary.
\item \textbf{The neuroscience is about acute stress.} Correct. The stress literature does not directly test whether weaker executive function predicts worse CBT outcome. Section~\ref{sec:moderators} reviews the sparse direct evidence.
\item \textbf{Therapists scaffold the executive load.} Plausible and untested. A therapist who generates the evidence list for the patient lowers the demand and could erase the moderator effect.
\end{enumerate}
\textbf{Limitations of this review.} It is a narrative critical review, written by a non-clinician; search was targeted rather than exhaustive in this draft; and the claims about mechanism are theoretical.
